\documentclass{cs0500}

\hwk{9}{Graph Algorithms - Part 2}{April 15, 2026 at 11:59 pm}

\begin{document}
\maketitle


\begin{problem}
    We study the relationship between maximum matchings and minimum vertex covers in bipartite graphs.
    \\ \textbf{Definition: } A vertex cover of an undirected graph $G = (V, E)$ is a subset of nodes $C \subseteq V$ such that every edge in $E$ has at least one endpoint in $C$. 
    \\ \textbf{König’s Theorem.} In any bipartite graph, the size of a minimum vertex cover equals the size of a maximum matching.

    \medskip 
    Let $G = (L \cup R, E)$ be a bipartite graph, where $L$ and $R$ are left and right nodes, respectively. Let $M^*$ be a maximum matching in $G$. Let $C^*$ be a minimum vertex cover of $G$. By part (a), $|C^*| \ge |M^*|$. To show $|C^*| \le |M^*|$, we construct a vertex cover $C$ of size $|M^*|$, and hence $|C^*| \le |C| = |M^*|$.
    \medskip 
    \\ Define a directed graph $G'=(V',E')$ with $V'=V\cup\{s,t\}$ as follows:
    \begin{itemize}
    \item $(s, u)$ for all unmatched $u \in L$.
    \item $(v, t)$ for all unmatched $v \in R$.
    \item $(u, v)$ for all $(u, v) \in (E \setminus M^*)$ with $u \in L$ and $v \in R$.
    \item $(v, u)$ for all $(u, v) \in M^*$ with $u \in L$ and $v \in R$.
    \end{itemize}

    There is a $s$-to-$t$ path in $G'$ if and only if there is an augmenting path with respect to $M^*$ in $G$.

    \textbf{Lemma 1. } Let $S$ be the set of nodes reachable from s in $G'$. Then $C = (L \ S) \cup (R \cap S)$ is a vertex cover of $G$ and $\vert C \vert = \vert M^* \vert$.
    \begin{enumerate}
        \item [(a)] Fix any undirected graph $G$. Let $M$ be a matching in $G$. Let $C$ be a vertex cover of $G$. Prove that $\vert C \vert \geq \vert M \vert$. (Hint: Can one node in $C$ cover multiple edges in $M$?)
        \item [(b)] Prove Lemma 1.
    \end{enumerate}
\end{problem}
\begin{solution}
    
\end{solution}

\newpage
\begin{problem}
Two departments jointly host a PhD visit day. Each of the $s$ students requests meetings with two faculty members: one from Department A and one from Department B. A faculty member may meet with any number of students. You must select a minimum-size set of faculty such that every student meets with at least one of their requested faculty members. The input consists of integers $f>0$ and $s>0$, where $f$ is the number of faculty in each department and $s$ is the number of students, followed by $s$ pairs $(a_i,b_i)$, where student $i$ requests faculty $a_i$ from Department A and faculty $b_i$ from Department B. Design an algorithm that outputs the minimum number of faculty required. Briefly describe your algorithm. Your algorithm should run in $O(f + s)$ time. 
\\ \emph{Hint: Use the bipartite graph construction and vertex-cover characterization from the previous problem.}

\end{problem}
\begin{solution}
    
\end{solution}
\newpage 
\begin{problem}
In class, we showed that a maximum matching in a bipartite graph with $n$ vertices
and $m$ edges can be computed in $O(mn)$ time. In this problem, we study the
\emph{Hopcroft--Karp algorithm}, which computes a maximum matching in
$O(m\sqrt{n})$ time.

The key idea of the algorithm is to repeatedly find \emph{shortest augmenting paths},
and more specifically, a maximal set of node-disjoint shortest augmenting paths.
The length of a path is defined as the number of edges it contains.

\medskip

\noindent
\textbf{Algorithm 1 (Hopcroft--Karp).}
\begin{itemize}
    \item \textbf{Input:} A bipartite graph $G$
    \item \textbf{Output:} A maximum matching $M$ of $G$
\end{itemize}

\begin{algorithmic}
\State $M \gets \emptyset$
\While{there exists an augmenting path in $G$ with respect to $M$}
    \State $\ell \gets$ the length of a shortest augmenting path
    \State $(P_1,\dots,P_k) \gets$ a maximal set of length-$\ell$ node-disjoint augmenting paths
    \State $M \gets M \oplus (P_1 \cup \cdots \cup P_k)$
\EndWhile
\State \Return $M$
\end{algorithmic}

\medskip

\noindent
\textbf{Lemma 2.} After each iteration of the \textbf{while} loop in Algorithm 1,
the length of the shortest augmenting path increases by at least $2$.

\medskip

\noindent
\textbf{Lemma 3.} Each iteration of the \textbf{while} loop in Algorithm 1
can be implemented to run in $O(m)$ time.

\medskip

\begin{enumerate}[(a)]
    \item Let $M$ be a matching in $G$, and let $M^*$ be a maximum matching.
    Suppose the shortest augmenting path with respect to $M$ has length $\ell$.
    Prove that
    \[
        \ell \cdot (|M^*| - |M|) \le n.
    \]
    \emph{Hint: Consider the graph $M \oplus M^*$. You may use any facts discussed
    in class without proof.}

    \item Using Lemma~2 and part (a), prove that the number of iterations of the
    \textbf{while} loop in Algorithm 1 is $O(\sqrt{n})$.

    \emph{Hint: Show that after $O(\sqrt{n})$ iterations, $|M|$ is close to
    $|M^*|$.}
\end{enumerate}
\end{problem}
\begin{solution}
    
\end{solution}



\end{document}
